% Self-contained section. Requires amsmath, amssymb, amsthm, hyperref.
% All labels/citation keys use h2: prefix. No custom macros required.
\section{Two independent exponents at a harmonic resonance}
\label{h2:sec:main}

The fully shifted height-one theorem in the incoming exact-identities
report varies the outer exponent and keeps every inner exponent equal
to one. Its final research section asks what survives when inner
exponents move independently. We answer a precise part of that question:
the first inner exponent may vary independently of the outer exponent,
with an arbitrary number of subsequent unit exponents. Every principal
Laurent coefficient and the finite part along every transverse affine
direction have a finite Gamma--Bernoulli formula. A polynomial in the
common shift supplies the same correction in all directions.

The meromorphic continuation of multiple zeta functions and the
dependence of specializations on the approach direction are classical
\cite{h2:AET,h2:MOW}. In particular, Matsumoto, Onozuka and Wakabayashi
give algorithms for Laurent expansions at arbitrary integer points,
allowing remainder integrals in general. The result below is a compact
closed coefficient formula for this particular common-shift family,
uniform in its trailing depth and its two-dimensional direction. We do
not claim priority for general continuation or multivariable Laurent
expansion. We also do not resolve arbitrary independent perturbations of
all trailing exponents.

\subsection{The family and its finite residue polynomial}

Throughout this section, $\Re a>0$, and powers of $n+a$ use the principal
logarithm. Define
\begin{align}
 e_r(n;a)&=[u^r]\prod_{j=0}^{n-1}\left(1+\frac{u}{j+a}\right),
 &p_n(a,u)&=\frac{(a+u)_n}{(a)_n},\label{h2:eq:e}\\
 Z_r(s,t;a)&=\sum_{n_0>n_1>\cdots>n_{r+1}\ge0}
 \frac{1}{(n_0+a)^s(n_1+a)^t\prod_{j=2}^{r+1}(n_j+a)},
 &&r\ge0.\label{h2:eq:Z}
\end{align}
The final product is empty at $r=0$. Thus $Z_0$ is a double Hurwitz
zeta function, and $Z_r(s,1;a)$ is the height-one function with $r+1$
inner unit exponents. Its depth generator is
\begin{equation}\label{h2:eq:generator}
 \mathcal Z_a(s,t,u)=\sum_{n\ge0}
     p_n(a,u)(n+a)^{-t}\zeta(s,n+a+1)
       =\sum_{r\ge0}Z_r(s,t;a)u^r.
\end{equation}
For example, absolute normal convergence holds when
$\Re s>1$ and $\Re(s+t)>2+\Re u$; a small closed disc in $u$
can be chosen in any strict version of these inequalities. The estimate
$p_n=O(n^{\Re u})$ and the usual Hurwitz-tail estimate prove this
claim and justify coefficient extraction. At each fixed depth,
$e_r(n;a)=O((1+\log n)^r)$ gives the corresponding normal convergence.

Let $B_k(x)$ denote Bernoulli polynomials, with $B_1(0)=-1/2$.
Define polynomials $q_j(u)$ through the formal series
\begin{equation}\label{h2:eq:qgen}
 \sum_{j\ge0}q_j(u)z^j
 =\exp\left(\sum_{k\ge1}
       \frac{(-1)^{k+1}\{B_{k+1}(u)-B_{k+1}(0)\}}
            {k(k+1)}z^k\right).
\end{equation}
This is a coefficient definition; no convergence of the infinite
asymptotic series is asserted. The classical Gamma-ratio expansion
\cite{DLMFGamma} reads
\begin{equation}\label{h2:eq:ratio}
 \frac{\Gamma(x+u)}{\Gamma(x)}
   \sim x^u\sum_{j\ge0}q_j(u)x^{-j}.
\end{equation}
In particular,
\[
 \begin{aligned}
 q_0&=1,& q_1&=\frac{u(u-1)}2,\\
 q_2&=\frac{u(u-1)(u-2)(3u-1)}{24},&
 q_3&=\frac{u^2(u-1)^2(u-2)(u-3)}{48}.
 \end{aligned}
\]
Every finite truncation has a locally uniform remainder in $u$.
Put
\begin{align}
 C_a(u)&=\frac{\Gamma(a)}{\Gamma(a+u)},\label{h2:eq:C}\\
 A_0(s)&=\frac1{s-1},\quad A_1(s)=-\frac12,\quad
 A_{2k}(s)=\frac{B_{2k}}{(2k)!}(s)_{2k-1}\ (k\ge1),\quad
 A_{2k+1}(s)=0\ (k\ge1),\label{h2:eq:A}\\
 P_m(s,u)&=\sum_{j=0}^{m+1}q_j(u)A_{m+1-j}(s).
 \label{h2:eq:P}
\end{align}
Thus $P_m$ is a finite rational expression, analytic for $s$ near $-m$.
Equivalently,
\[
 P_m(s,u)=\frac{q_{m+1}(u)}{s-1}-\frac{q_m(u)}2
  +\sum_{k=1}^{\lfloor(m+1)/2\rfloor}
       \frac{B_{2k}(s)_{2k-1}}{(2k)!}q_{m+1-2k}(u).
\]
The first two examples are
\begin{align}
 P_0(s,u)&=\frac{u(u-1)}{2(s-1)}-\frac12,
 \label{h2:eq:P0}\\
 P_1(s,u)&=\frac{u(u-1)(u-2)(3u-1)}{24(s-1)}
             -\frac{u(u-1)}4+\frac{s}{12}.
 \label{h2:eq:P1}
\end{align}

\begin{proposition}[One moving divisor before depth extraction]
\label{h2:prop:local}
For every integer $m\ge0$, there is a jointly holomorphic function
$\mathcal H_{m,a}$ near $(s,t,u)=(-m,1,0)$ such that
\begin{equation}\label{h2:eq:local}
 \mathcal Z_a(s,t,u)=
 \frac{C_a(u)P_m(s,u)}{s+t+m-1-u}+\mathcal H_{m,a}(s,t,u).
\end{equation}
The statement is also locally uniform in $a$ on the right half-plane.
\end{proposition}
\begin{proof}
Write $x=n+a$. Since
$\zeta(s,x+1)=\zeta(s,x)-x^{-s}$, the Hurwitz asymptotic expansion
\cite{DLMFHzeta} gives
\[
 \zeta(s,x+1)\sim\sum_{k\ge0}A_k(s)x^{1-s-k}.
\]
Multiplying a finite truncation by \eqref{h2:eq:ratio} yields
\[
 p_n(a,u)x^{-t}\zeta(s,x+1)
 \sim C_a(u)\sum_{\ell\ge0}
   \left(\sum_{j+k=\ell}q_j(u)A_k(s)\right)
                         x^{1-s-t+u-\ell}.
\]
Subtract sufficiently many terms from the summand in
\eqref{h2:eq:generator}. The remainder series converges normally in a
neighborhood of $(-m,1,0)$: after retaining $L$ terms, its
summands satisfy
\[
 O(n^{1-\Re(s+t)+\Re u-L}).
\]
Choosing $L>m+2$ suffices on a small neighborhood. Parameter derivatives introduce only fixed powers
of $\log n$ after shrinking that neighborhood.

The sums of the subtracted terms are the corresponding finite linear
combination of
$\zeta(s+t-1-u+\ell,a)$. Exactly one can have a pole in this
neighborhood: $\ell=m+1$. Its residue is
$C_a(u)P_m(s,u)$, and subtracting its simple pole leaves a holomorphic
function. This proves \eqref{h2:eq:local}. The use of finite asymptotic
subtraction, with a normally convergent remainder, is what makes this a
meromorphic identity rather than a formal asymptotic assertion.
\end{proof}

\subsection{An explicit value of the holomorphic remainder}

The next lemma is the specialization mechanism from the incoming
fully shifted height-one theorem \cite{IncomingExact}. We include its
proof because the direction-independent correction in the main theorem
depends on it.

\begin{lemma}[A finite Stirling specialization]
\label{h2:lem:stirling}
Let
\[
 F_a(w,u)=\sum_{n\ge0}\frac{p_n(a,u)}{(n+a)^w},\qquad
 T_m(a,u)=-\sum_{j=1}^{m+1}
   \left\{\begin{matrix}m+1\\j\end{matrix}\right\}
      \frac{(a-1)^{\underline j}}{u+j}.
\]
Here $(a-1)^{\underline j}$ is a falling factorial and braces denote
Stirling numbers of the second kind. For sufficiently small nonzero
$u$, the meromorphic continuation satisfies
\begin{equation}\label{h2:eq:Fnegative}
 F_a(-m,u)=T_m(a,u).
\end{equation}
Consequently, with
\begin{equation}\label{h2:eq:L}
 L_m(a,u)=\frac{T_m(a,u)-\zeta(-m,a)}u,
\end{equation}
one has
\begin{equation}\label{h2:eq:Hvalue}
 \mathcal H_{m,a}(-m,1,u)
   =L_m(a,u)+\frac{C_a(u)P_m(-m,u)}u.
\end{equation}
The right side of \eqref{h2:eq:Hvalue} is holomorphic at $u=0$.
\end{lemma}
\begin{proof}
The Mellin kernel of $F_a$ is
\[
 G_a(v,u)=e^{-av}\,{}_2F_1(1,a+u;a;e^{-v}).
\]
Gauss connection at one gives
\begin{align*}
 G_a(v,u)&=\frac{\Gamma(a)\Gamma(1+u)}{\Gamma(a+u)}
       \frac{e^{-v}}{(1-e^{-v})^{1+u}}+R_a(v,u),\\
 R_a(v,u)&=-\frac{a-1}{1+u}e^{-av}
      {}_2F_1(1,a+u;2+u;1-e^{-v}).
\end{align*}
The second kernel is analytic at $v=0$ and decays exponentially at
infinity, locally in the parameters. Mellin subtraction therefore
shows that its quotient by $\Gamma(w)$ is entire in $w$, and has value
$(-1)^m m![v^m]R_a(v,u)$ at $w=-m$. The first kernel has local
exponents $-1-u+j$; for small nonzero $u$ its Mellin transform has no
pole at $w=-m$, so division by $\Gamma(w)$ makes its value zero there.
This proves that the left side of \eqref{h2:eq:Fnegative} is a
polynomial in $a$ of degree at most $m+1$.

To identify the polynomial, take a positive integer $a=A$. Direct
index shifting in the original Dirichlet series gives
\[
 F_A(w,u)=\frac{\Gamma(A)\Gamma(1+u)}{\Gamma(A+u)}
 \left(F_1(w,u)-\sum_{k=1}^{A-1}
       \frac{(1+u)_{k-1}}{(k-1)!k^w}\right).
\]
At $w=-m$, the first term vanishes. The finite identity
\[
 k^m=\sum_{j=1}^{m+1}
    \left\{\begin{matrix}m+1\\j\end{matrix}\right\}
                           (k-1)^{\underline{j-1}}
\]
and the binomial summation
\[
 \frac{\Gamma(A)\Gamma(1+u)}{\Gamma(A+u)}
 \sum_{k=1}^{A-1}\frac{(1+u)_{k-1}}{(k-1)!}
           (k-1)^{\underline{j-1}}
       =\frac{(A-1)^{\underline j}}{u+j}
\]
give $T_m(A,u)$. Both sides are polynomials in $a$ of degree at most
$m+1$, and agreement at all positive integers proves
\eqref{h2:eq:Fnegative} for every $\Re a>0$.

Finally, $p_{n+1}-p_n=up_n/(n+a)$ implies, first by absolutely
convergent summation,
\begin{equation}\label{h2:eq:telescoping}
 \mathcal Z_a(s,1,u)=\frac{F_a(s,u)-\zeta(s,a)}u.
\end{equation}
Meromorphic continuation, \eqref{h2:eq:local}, and
\eqref{h2:eq:Fnegative} now give \eqref{h2:eq:Hvalue}. Its
holomorphy also follows directly from Proposition~\ref{h2:prop:local}.
\end{proof}

The order of operations in \eqref{h2:eq:Fnegative} matters. It restricts
$w$ to $-m$ for nonzero $u$ and then continues the resulting function in
$u$. At a crossing of a moving pole it does not give the Laurent finite
part obtained by first extracting a coefficient in $u$. The main
theorem keeps these operations in the required order.

For later use put
\begin{equation}\label{h2:eq:Delta}
 \Delta_{m,d}(a)=(-1)^{d+1}\sum_{j=1}^{m+1}
  \left\{\begin{matrix}m+1\\j\end{matrix}\right\}
                      \frac{(a-1)^{\underline j}}{j^{d+1}}.
\end{equation}
Thus $[u^r]L_m(a,u)=\Delta_{m,r+1}(a)$ for $r\ge0$.

\subsection{The complete two-direction finite-part theorem}

\begin{theorem}[All depths and every transverse affine direction]
\label{h2:thm:direction}
Let $m,r\ge0$, $\Re a>0$, and $\alpha,\beta\in\mathbb C$ with
$q=\alpha+\beta\ne0$. Set $\theta=\alpha/q$ and define the analytic
germ
\begin{equation}\label{h2:eq:V}
 V_{m,\theta}(u;a)=C_a(u)P_m(-m+\theta u,u).
\end{equation}
Then
\begin{equation}\label{h2:eq:main}
 \boxed{\displaystyle
 Z_r(-m+\alpha\varepsilon,1+\beta\varepsilon;a)
 =\sum_{h=0}^{r+1}
     \frac{[u^{r+1-h}]V_{m,\theta}(u;a)}
          {q^h\varepsilon^h}
      +\Delta_{m,r+1}(a)+O(\varepsilon).}
\end{equation}
In particular,
\begin{equation}\label{h2:eq:FP}
 \operatorname{FP}_{\varepsilon=0}
 Z_r(-m+\alpha\varepsilon,1+\beta\varepsilon;a)
 =[u^{r+1}]V_{m,\theta}(u;a)+\Delta_{m,r+1}(a).
\end{equation}
Every displayed coefficient is a finite expression in rational numbers,
$a$, $\theta$, $q^{-1}$, and the polygamma values
\[
 \psi(a),\psi'(a),\ldots,\psi^{(r)}(a).
\]
\end{theorem}
\begin{proof}
In \eqref{h2:eq:local}, substitute the affine path and put
$f(\varepsilon,u)=C_a(u)P_m(-m+\alpha\varepsilon,u)$. The singular
term becomes $f(\varepsilon,u)/(q\varepsilon-u)$. For fixed nonzero
$\varepsilon$, expand in $u$ first:
\[
 \frac1{q\varepsilon-u}
    =\sum_{n\ge0}\frac{u^n}{q^{n+1}\varepsilon^{n+1}}.
\]
If $h\ge1$, collecting the coefficient of
$u^r\varepsilon^{-h}$ gives
\[
 q^{-h}[u^{r+1-h}]
       C_a(u)P_m\left(-m+\frac{\alpha}{q}u,u\right).
\]
For the constant Laurent coefficient the same calculation excludes the
term of degree zero in $\varepsilon$ in $f$, and therefore gives
\[
 [u^{r+1}]C_a(u)
   \left(P_m(-m+\theta u,u)-P_m(-m,u)\right).
\]
The holomorphic remainder contributes
$[u^r]\mathcal H_{m,a}(-m,1,u)$ to the constant coefficient and
nothing to the principal part. Formula \eqref{h2:eq:Hvalue} cancels
$-[u^{r+1}]C_a(u)P_m(-m,u)$ and leaves
$[u^r]L_m=\Delta_{m,r+1}$. This proves \eqref{h2:eq:main}.

Finally,
\[
 C_a(u)=\exp\left(-\sum_{k\ge1}
                  \frac{\psi^{(k-1)}(a)}{k!}u^k\right).
\]
Only finitely many coefficients, through degree $r+1$, enter the
formula. This proves the stated finite algebra and completes the proof.
\end{proof}

The choice $(\alpha,\beta)=(1,0)$ recovers the archived common-shift
height-one expansion. Indeed comparison of the moving residues in
\eqref{h2:eq:telescoping} proves the finite identity
\begin{equation}\label{h2:eq:recovery}
 uP_m(-m+u,u)=q_{m+1}(u).
\end{equation}
If the former notation is used, $q_{m+1}$ is precisely its polynomial
$W_m$. The choice $(\alpha,\beta)=(0,1)$ instead keeps the outer
exponent at $-m$ and moves the first inner exponent. Formula
\eqref{h2:eq:FP} specifies their discrepancy at every depth; it is
not legitimate to transfer a finite part from one direction to the
other without that correction.

\begin{corollary}[Pole order is independent of the transverse direction]
\label{h2:cor:poles}
For $m=0$ or positive odd $m$, the pole in \eqref{h2:eq:main} has order
exactly $r+1$, with leading coefficient
$\zeta(-m)/q^{r+1}$. For positive even $m$, it has order exactly $r$
when $r\ge1$, with leading coefficient
\[
 \frac{m+1}{2m}\frac{B_m}{q^r}.
\]
At positive even $m$ and $r=0$ the function is holomorphic at the
intersection, and the value is independent of the direction.
\end{corollary}
\begin{proof}
From \eqref{h2:eq:qgen}, $q_0(0)=1$ and $q_j(0)=0$ for $j>0$.
Hence $P_m(-m,0)=A_{m+1}(-m)=\zeta(-m)$. This proves the first
assertion. For positive even $m$, $A_{m+1}$ is identically zero, so
$[u]P_m(-m+\theta u,u)$ is independent of $\theta$. It can be
computed at $\theta=1$ using \eqref{h2:eq:recovery}.

For completeness, the Gamma-ratio coefficients also obey
\[
 q_{m+1}(u)=(u-m)_{m+1}[v^{m+1}]
                e^{-v}\left(\frac{v}{1-e^{-v}}\right)^{1+u}.
\]
To justify this auxiliary identity, the beta integral gives, for
$\Re u<0$ and sufficiently large positive $x$,
\[
 \frac{\Gamma(x+u)}{\Gamma(x)}
 =\frac1{\Gamma(-u)}\int_0^\infty e^{-xv}v^{-1-u}
       e^{-uv}\left(\frac{v}{1-e^{-v}}\right)^{1+u}\,dv.
\]
Expanding the analytic factor at $v=0$ and integrating a finite Taylor
polynomial gives
\[
 q_j(u)=(-u)_j[v^j]\,
          e^{-uv}\left(\frac{v}{1-e^{-v}}\right)^{1+u}.
\]
The analytic factor here is obtained from
$e^{-v}(v/(1-e^{-v}))^{1+u}$ by replacing $v$ with $-v$.
Since $(-1)^j(-u)_j=(u-j+1)_j$, this proves the identity with
$j=m+1$. Both sides are polynomials in $u$, so it holds for every $u$
by polynomial continuation.
For even $m\ge2$, the coefficient at $u=0$ in the last bracket
vanishes, while its coefficient of $u$ is
\[
 [v^{m+1}]\frac{v}{e^v-1}
                   \log\frac{v}{1-e^{-v}}
       =\frac{m+1}{2m}\frac{B_m}{m!}.
\]
This follows on multiplying the ordinary Bernoulli series by
$v/2-\sum_{k\ge1}B_{2k}v^{2k}/(2k(2k)!)$.
Since $(u-m)_{m+1}=m!u+O(u^2)$, one obtains
$[u^2]q_{m+1}=(m+1)B_m/(2m)\ne0$. Formula
\eqref{h2:eq:main} now gives the stated order and leading coefficient.
For $r=0$, the local numerator $P_m(s,0)=A_{m+1}(s)$ vanishes
identically, proving joint holomorphy.
\end{proof}

\paragraph{Directions within the polar hyperplane.}
The condition $q\ne0$ is essential whenever the pole in
Corollary~\ref{h2:cor:poles} is genuine. A nonconstant affine path with
$\alpha+\beta=0$ remains inside $s+t+m-1=0$, so the meromorphic function
has no ordinary one-variable restriction along that path. The exception
is $r=0$ and positive even $m$, where the function is jointly holomorphic:
all approach directions, including tangent ones, give
\[
 Z_0(-m,1;a)=\Delta_{m,1}(a)+\frac{m+1}{2m}B_m.
\]

\subsection{Explicit identities and direction corrections}

Write $P=\psi(a)$ and $Q=\psi'(a)$. Formula \eqref{h2:eq:P0}
gives
\[
 P_0(\theta u,u)=-\frac12+\frac u2
       +\frac{\theta-1}{2}u^2
       +\frac{\theta^2-\theta}{2}u^3+O(u^4).
\]
Consequently,
\begin{align}
 Z_0(\alpha\varepsilon,1+\beta\varepsilon;a)
 &= -\frac1{2q\varepsilon}+a-\frac12+\frac P2
                   +O(\varepsilon),\label{h2:eq:example00}\\
 Z_1(\alpha\varepsilon,1+\beta\varepsilon;a)
 &= -\frac1{2q^2\varepsilon^2}
      +\frac{1+P}{2q\varepsilon}
      +1-a+\frac{\theta-1}{2}-\frac P2
             -\frac{P^2-Q}{4}+O(\varepsilon).
 \label{h2:eq:example01}
\end{align}
Thus at total depth three the finite part changes by exactly
$(\theta-\eta)/2$ when the direction parameter changes from $\eta$
to $\theta$; the result holds at every common shift $a$.

At the next outer integer one obtains
\begin{equation}\label{h2:eq:example10}
 Z_0(-1+\alpha\varepsilon,1+\beta\varepsilon;a)
 =-\frac1{12q\varepsilon}
    +\frac{a^2+a}{4}+\frac{2\theta-5}{24}+\frac P{12}
                  +O(\varepsilon).
\end{equation}
The two coordinate directions differ by $1/12$ in their finite parts.
For an example where the highest potential pole cancels, put
$L=\gamma+2\log2$. At the half shift,
\begin{equation}\label{h2:eq:example21}
 Z_1(-2+\alpha\varepsilon,1+\beta\varepsilon;1/2)
    =\frac1{8q\varepsilon}+\frac L8-\frac{7+12\theta}{288}
                           +O(\varepsilon).
\end{equation}
At $\theta=1$ this is the previously obtained height-one finite part
$L/8-19/288$; at $\theta=0$ it is $L/8-7/288$.

In general define the difference between two directional finite parts
by $\mathfrak A_{m,r}(\theta,\eta;a)$. It has the exact form
\begin{equation}\label{h2:eq:anomaly}
 \mathfrak A_{m,r}(\theta,\eta;a)
 =[u^{r+1}]C_a(u)
       \bigl(P_m(-m+\theta u,u)-P_m(-m+\eta u,u)\bigr).
\end{equation}
The Stirling polynomial has canceled completely. At double-zeta depth
$r=0$ this reduces to
\begin{equation}\label{h2:eq:anomaly0}
 \mathfrak A_{m,0}(\theta,\eta;a)
   =-(\theta-\eta)\zeta(-m)H_m,
 \qquad H_0=0.
\end{equation}
To see the last equality, $[u]$ of the difference is
$(\theta-\eta)A'_{m+1}(-m)$. For odd $m$,
$A_{m+1}(s)=B_{m+1}(s)_m/(m+1)!$, and its logarithmic derivative
at $s=-m$ is $-H_m$; the cases $m=0$ and positive even $m$ vanish.
Thus the double-zeta directional correction is rational and independent
of $a$, even though each individual finite part contains $\psi(a)$.

\subsection{An exact normal slice and all Stieltjes derivatives}

At depth two one can go beyond the finite part in the direction that
keeps the outer exponent fixed. The resulting reduction is a classical
Bernoulli specialization of a multiple zeta function; it also supplies
particularly short mixed Gamma--Stieltjes identities.

\begin{theorem}[The complete inner-exponent slice]
\label{h2:thm:slice}
For $m\ge0$, the meromorphic identity
\begin{equation}\label{h2:eq:slice}
 Z_0(-m,t;a)=-\frac1{m+1}\sum_{j=0}^{m+1}
       \binom{m+1}{j}B_j(1)\zeta(t-m-1+j,a)
\end{equation}
holds. More generally, for every trailing depth $r$,
\begin{equation}\label{h2:eq:slicedepth}
 Z_r(-m,t;a)=-\frac1{m+1}\sum_{j=0}^{m+1}
       \binom{m+1}{j}B_j(1)M_r(t-m-1+j;a),
\end{equation}
where $M_r(w;a)=\sum_{n\ge0}e_r(n;a)(n+a)^{-w}$.
\end{theorem}
\begin{proof}
In the half-plane $\Re t>m+2$, the sum
$\sum e_r(n;a)(n+a)^{-t}\zeta(s,n+a+1)$ is locally normally
convergent for $s$ near $-m$. The Bernoulli evaluation
\[
 \zeta(-m,x+1)=-\frac{B_{m+1}(x+1)}{m+1}
 =-\frac1{m+1}\sum_{j=0}^{m+1}
       \binom{m+1}{j}B_j(1)x^{m+1-j}
\]
gives the finite formula termwise there. Continuing meromorphically in
$t$ proves the identities. This argument specifies the restriction
$s=-m$ before taking a Laurent coefficient in $t$.
\end{proof}

Use the Stieltjes convention
\[
 \zeta(1+\varepsilon,a)=\frac1\varepsilon
       +\sum_{\ell\ge0}\frac{(-1)^\ell}{\ell!}
                         \gamma_\ell(a)\varepsilon^\ell,
 \qquad \gamma_0(a)=-\psi(a).
\]
Writing $\zeta^{(\ell)}(w,a)=\partial_w^\ell\zeta(w,a)$,
\eqref{h2:eq:slice} yields the complete regular coefficients
\begin{align}
 [\varepsilon^\ell]Z_0(-m,1+\varepsilon;a)
 ={}&-\frac1{(m+1)\ell!}\sum_{j=0}^{m}
       \binom{m+1}{j}B_j(1)\zeta^{(\ell)}(j-m,a)
       \notag\\
 &+\frac{(-1)^\ell}{\ell!}\zeta(-m)\gamma_\ell(a),
 \qquad \ell\ge0.\label{h2:eq:slicejets}
\end{align}
Its principal part is $\zeta(-m)/\varepsilon$. At $m=0$ this simplifies to
\begin{equation}\label{h2:eq:allstieltjes}
 \boxed{\displaystyle
 \operatorname{FP}_{t=1}\partial_t^\ell Z_0(0,t;a)
   =-\zeta^{(\ell)}(0,a)
           +\frac{(-1)^{\ell+1}}2\gamma_\ell(a),
 \qquad \ell\ge0.}
\end{equation}
For $\ell=1$, Lerch's classical identity
\cite{DLMFHzeta} gives the concrete logarithmic formula
\begin{equation}\label{h2:eq:firststieltjes}
 \boxed{\displaystyle
 \operatorname{FP}_{t=1}\partial_t Z_0(0,t;a)
   =-\log\Gamma(a)+\frac12\log(2\pi)+\frac12\gamma_1(a).}
\end{equation}
The logarithm of $\Gamma$ is its analytic branch on $\Re a>0$,
real for positive $a$. Equivalently,
\begin{align}
 Z_0(0,1+\varepsilon;a)
  ={}&-\frac1{2\varepsilon}+a-\frac12+\frac12\psi(a)\notag\\
  &+\varepsilon\left(-\log\Gamma(a)+\frac12\log(2\pi)
                                     +\frac12\gamma_1(a)\right)
       +O(\varepsilon^2).\label{h2:eq:Z0next}
\end{align}
This does not evaluate the first regular coefficient in the different
slice $Z_0(\varepsilon,1;a)$; the two slices coincide only at those
coefficients explicitly identified above.

There is also a finite antiderivative at every Stieltjes order. Put
$J_\ell(a)=\operatorname{FP}_{t=1}\partial_t^\ell Z_0(0,t;a)$.
Then
\begin{equation}\label{h2:eq:primitive}
 \boxed{\displaystyle
 \int J_\ell(a)\,da
  =\frac{\zeta^{(\ell+1)}(0,a)}{2(\ell+1)}
      -\ell!\sum_{j=0}^{\ell}
                  \frac{\zeta^{(j)}(-1,a)}{j!}
       +\text{constant}.}
\end{equation}
Indeed the Hurwitz shift derivative implies
\[
 \partial_a\zeta^{(j)}(-1,a)
     =\zeta^{(j)}(0,a)-j\zeta^{(j-1)}(0,a),\qquad
 \partial_a\zeta^{(\ell+1)}(0,a)
     =(-1)^{\ell+1}(\ell+1)\gamma_\ell(a).
\]
The finite sum telescopes on differentiation. For example, a primitive
of the right side of \eqref{h2:eq:firststieltjes} is
\[
 \frac14\zeta''(0,a)-\zeta'(-1,a)-\zeta(-1,a).
\]
These formulas provide exact derivative and antiderivative identities
in the independent inner spectral direction, including every
generalized Stieltjes order.

\subsection{Scope, verification, and the next questions}

The new finite-part theorem stops at $O(\varepsilon)$ because the
first derivatives of $\mathcal H_{m,a}$ can enter that term.
The normal slice at depth two is exceptional: the Bernoulli reduction
evaluates all of its regular coefficients. Formula
\eqref{h2:eq:slicedepth} reduces the corresponding problem at larger
depth by one, but does not by itself reduce every resulting spectral
derivative to ordinary Hurwitz data.

The supplied program \texttt{verify\_two\_exponent.py} checks the
residue identity \eqref{h2:eq:recovery} exactly for $0\le m\le6$.
It separately evaluates the continued nested sum using its original
finite initial segment and a long Euler--Maclaurin tail. Discrete Cauchy
averages then extract Laurent coefficients without using the predicted
principal part in the extraction. A second cutoff and longer tail
check numerical stability. Eight cases at $65$ decimal working digits
compare $21$ Laurent coefficients through the finite part; the maximum
absolute discrepancy is below $2.25\times10^{-46}$. The longer-tail
repeat differs by less than $5.9\times10^{-52}$. These are
arbitrary-precision diagnostics;
the normal-convergence and coefficient-extraction proofs above establish
the identities for all parameters.

Three further research directions are concrete. First, allow the
second inner exponent to move as well; this introduces several
intersecting polar divisors and tests whether a finite analogue of
$P_m$ still eliminates every constant-term remainder. Second, determine
linear combinations of first regular coefficients in different
transverse directions that cancel the derivatives of
$\mathcal H_{m,a}$. Third, compare the rational directional correction
at $a=1$ with the general Laurent algorithms and with the
Stirling-polynomial formulas for nonpositive multiple-zeta values
\cite{h2:MOW,h2:IshiiShinohara}; this requires matching both the order of
indices and the chosen regularization, since the present base point has
unit inner exponents.
